\documentclass[12pt,a4paper]{report}
\usepackage[utf-8]{inputenc}
\usepackage[french]{babel}
\usepackage[margin=2cm]{geometry}
\usepackage{graphicx}
\usepackage{xcolor}
\usepackage{listings}
\usepackage{fancyhdr}
\usepackage{hyperref}
\usepackage{array}
\usepackage{booktabs}
\usepackage{amssymb}
\usepackage{tikz}
\usepackage{float}

% Configuration Listings pour Rust
\lstset{
  language=Rust,
  basicstyle=\ttfamily\small,
  breaklines=true,
  keywordstyle=\color{blue}\bfseries,
  commentstyle=\color{gray},
  stringstyle=\color{red},
  showstringspaces=false,
  tabsize=2,
  frame=single,
  backgroundcolor=\color{lightgray!20},
  numbers=left,
  numberstyle=\tiny,
  stepnumber=1
}

\title{\textbf{QSDID Platform}\\
  \Large Gestion d'Identité et Encapsulation des Clés\\
  \normalsize Plateforme d'Identité Quantique Sécurisée Décentralisée}
\author{Architecture Cryptographique Post-Quantique}
\date{\today}

\pagestyle{fancy}
\fancyhf{}
\rhead{QSDID Platform}
\lhead{Documentation Technique}
\cfoot{\thepage}

\begin{document}

\maketitle

\tableofcontents
\newpage

% ============================================================================
\chapter{Vue d'Ensemble et Architecture}
% ============================================================================

\section{Vue d'Ensemble du Projet}

Le projet QSDID est une \textbf{Plateforme d'Identité Quantique Sécurisée Décentralisée} construite en \textbf{Rust}, inspirée par l'architecture IOTA Identity, mais fortifiée par la \textbf{cryptographie post-quantique (PQC)} et les signatures \textbf{hybrides} pour la résistance quantique.

\subsection{Caractéristiques Principales}

\begin{itemize}
  \item \textbf{13 crates Rust} spécialisées dans la gestion cryptographique
  \item Support de \textbf{ML-DSA-65} (signatures post-quantiques) + \textbf{Ed25519} (signatures classiques)
  \item Intégration \textbf{JOSE/JWK} complète pour la sérialisation
  \item Stockage distribué via \textbf{IPFS/BlockFrost}
  \item Conformité aux standards NIST FIPS 204 et 203
\end{itemize}

\subsection{Architecture Modulaire}

Le projet est organisé en 13 modules complémentaires :

\begin{enumerate}
  \item \texttt{identity\_core} : Types fondamentaux et utilitaires
  \item \texttt{identity\_did} : Implémentation des DIDs (Decentralized Identifiers)
  \item \texttt{identity\_credential} : Gestion des credentials
  \item \texttt{identity\_document} : Documents d'identité
  \item \texttt{identity\_jose} : Implémentation JOSE (JWK, JWS, JWT)
  \item \texttt{identity\_verification} : Vérification de signatures
  \item \texttt{identity\_pqc\_verifier} : Vérification post-quantique
  \item \texttt{pqc\_ml\_dsa} : Signatures ML-DSA-65 (FIPS 204)
  \item \texttt{pqc\_ml\_kem} : Encapsulation ML-KEM-768 (FIPS 203)
  \item \texttt{hybrid\_signer} : Signatures hybrides avec WNS
  \item \texttt{identity\_storage} : Stockage distribué et gestion blockchain
  \item \texttt{storage\_client} : Client de stockage
  \item \texttt{test\_pqc} : Tests de cryptographie post-quantique
\end{enumerate}

% ============================================================================
\chapter{Couche Gestion d'Identités}
% ============================================================================

\section{Architecture DID (Decentralized Identifier)}

\subsection{Concept Fondamental}

Un \textbf{DID} (Decentralized Identifier) est un identificateur unique, autogéré et vérifiable conforme aux standards W3C. Contrairement aux systèmes centralisés, les DIDs ne dépendent d'aucune autorité centrale.

\subsection{Structure DID}

\begin{lstlisting}
Format: did:method:network:identifier

Exemple: did:iota:main:axFkJEKz5a4
         |  |    |    |
         |  |    |    +-- Method-specific ID
         |  |    +------- Network (optional)
         |  +------------ Method name  
         +-------------- Scheme (always "did")
\end{lstlisting}

\subsection{Trait DID}

Le trait \texttt{DID} définit l'interface universelle pour tous les types d'identifiants :

\begin{lstlisting}[language=Rust]
pub trait DID: Clone + PartialEq + Eq + PartialOrd + Ord + Hash {
  const SCHEME: &'static str = "did";
  
  fn scheme(&self) -> &'static str;      // "did"
  fn method(&self) -> &str;               // "iota", "example"
  fn method_id(&self) -> &str;           // Identifiant unique
  fn authority(&self) -> &str;           // method:method_id
  fn as_str(&self) -> &str;              // Représentation sérialisée
  fn to_url(&self) -> DIDUrl;
}
\end{lstlisting}

\subsection{CoreDID}

La structure \texttt{CoreDID} implémente le trait \texttt{DID} :

\begin{itemize}
  \item \textbf{Parse et validation} stricte selon W3C standards
  \item \textbf{Méthodes setter} pour modification sécurisée
  \item \textbf{Validation} des caractères de method-name et method-id
  \item \textbf{Sérialisation} JSON complète
\end{itemize}

Exemple de validation :
\begin{lstlisting}[language=Rust]
let did = CoreDID::parse("did:iota:main:axFkJEKz5a4")?;

assert_eq!(did.scheme(), "did");
assert_eq!(did.method(), "iota");
assert_eq!(did.method_id(), "main:axFkJEKz5a4");
\end{lstlisting}

\section{DIDs Étendus}

\subsection{DID-JWK}

Encode directement une clé publique JWK dans l'identifiant :
\begin{lstlisting}
did:jwk:eyJjcnYiOiJFZDI1NTE5IiwieCI6IlZ...
\end{lstlisting}

Avantages :
\begin{itemize}
  \item Autonome (pas de blockchain)
  \item Vérification directe de la clé
  \item Résilience maximale
\end{itemize}

\subsection{DID-CompositeJWK}

Combine clés post-quantiques + classiques dans l'identifiant.

\section{Système de Credentials}

\subsection{Structure CredentialData}

\begin{lstlisting}[language=Rust]
#[derive(Serialize, Deserialize, Clone)]
pub struct CredentialData {
    pub claims: serde_json::Value,
    pub metadata: Metadata,
    pub ai_result: AiResult,
    pub image: Option<Vec<u8>>,
}

pub struct Metadata {
    pub credential_type: String,
    pub issuer_did: String,
    pub holder_did: String,
    pub uuid: String,
    pub issued_at: String,
    pub expires_at: Option<String>,
}

pub struct AiResult {
    pub fraud_score: f64,           // 0.0 to 1.0
    pub heatmap_hash: String,
    pub analysis_timestamp: String,
}
\end{lstlisting}

\subsection{Cycle de Vie du Credential}

\begin{enumerate}
  \item \textbf{Génération} : Émetteur crée le credential signé
  \item \textbf{Signature} : Composite signature (ML-DSA-65 + Ed25519)
  \item \textbf{Stockage} : Sauvegarde sur blockchain/IPFS via StorageClient
  \item \textbf{Distribution} : Holder reçoit rootCid
  \item \textbf{Vérification} : Holder récupère et valide la signature
  \item \textbf{Révocation} : Émetteur peut révoquer via credential registry
\end{enumerate}

% ============================================================================
\chapter{Couche Encapsulation des Clés}
% ============================================================================

\section{Architecture JOSE}

\textbf{JOSE} (JSON Object Signing and Encryption) est un framework standardisé pour :
\begin{itemize}
  \item Signature de données JSON (\textbf{JWS})
  \item Chiffrement de données JSON (\textbf{JWE})
  \item Tokens web (\textbf{JWT})
  \item Représentation de clés (\textbf{JWK})
\end{itemize}

\subsection{Structure des Modules}

\begin{lstlisting}
identity_jose/src/
├── jwk/                    # JSON Web Keys
│   ├── key.rs             # Clé générale
│   ├── key_hybrid.rs      # Clés hybrides
│   ├── composite_jwk.rs   # Clés composites PQC
│   ├── key_type.rs        # Types: EC, RSA, OKP, Akp
│   └── key_params.rs      # Paramètres
├── jws/                    # JSON Web Signature
│   ├── algorithm.rs       # Algorithmes JWS
│   └── encoding/          # Encodage/Décodage
├── jwt/                    # JSON Web Token
│   ├── claims.rs
│   └── header.rs
└── jose/traits.rs         # Traits JOSE
\end{lstlisting}

\section{JWK (JSON Web Key)}

\subsection{Types de Clés Supportées}

\begin{table}[H]
\centering
\begin{tabular}{|c|c|c|c|}
\hline
\textbf{Type} & \textbf{Notation} & \textbf{Usage} & \textbf{Taille} \\
\hline
RSA & \texttt{kty: "RSA"} & Chiffrement classique & 2048-4096 bits \\
\hline
EC (P-256/P-384) & \texttt{kty: "EC"} & ECDSA signatures & 256-384 bits \\
\hline
OKP (Ed25519) & \texttt{kty: "OKP"} & EdDSA signatures & 256 bits \\
\hline
AKP & \texttt{kty: "Akp"} & \textbf{Post-Quantum (ML-DSA-65)} & 1952 bytes \\
\hline
\end{tabular}
\caption{Types de clés JWK supportées}
\end{table}

\section{TraditionalJwk (Clés Classiques)}

\subsection{Structure}

\begin{lstlisting}[language=Rust]
pub struct TraditionalJwk {
    pub kty: JwkType,           // "RSA", "EC", "OKP"
    pub use_: Option<JwkUse>,  // "sig" ou "enc"
    pub key_ops: Vec<JwkOperation>,  // "sign", "verify", etc.
    pub alg: Option<String>,    // "ES256", "EdDSA", "RS256"
    pub kid: Option<String>,    // Key ID (thumbprint)
    pub x5u: Option<Url>,       // X.509 certificate URL
    pub x5c: Vec<String>,       // Certificate chain
    pub params: JwkParams,      // Paramètres spécifiques
}
\end{lstlisting}

\subsection{Paramètres OKP (Ed25519)}

\begin{lstlisting}[language=json]
{
  "kty": "OKP",
  "crv": "Ed25519",
  "x": "<base64url of public key>",
  "d": "<base64url of secret key>"
}
\end{lstlisting}

\begin{itemize}
  \item \textbf{x} : Clé publique (base64url)
  \item \textbf{d} : Clé privée (base64url, seulement si clé privée)
  \item \textbf{crv} : Courbe (Ed25519)
\end{itemize}

\section{PostQuantumJwk (Clés Post-Quantiques)}

\subsection{Structure}

\begin{lstlisting}[language=Rust]
pub struct PostQuantumJwk(Jwk) {
    pub kty: JwkType::Akp,      // Always "Akp"
    pub crv: "ML-DSA-65",       // Post-quantum curve
    pub x: Vec<u8>,             // Public key (1952 bytes)
    pub d: Option<Vec<u8>>,     // Secret key (2560 bytes)
    pub kid: String,            // Key identifier
}
\end{lstlisting}

\subsection{Restrictions et Validations}

\begin{itemize}
  \item \textbf{jty obligatoirement "Akp"} : Ne peut contenir que des clés post-quantiques
  \item \textbf{Algorithmes valides} : `id-MLDSA44-Ed25519` ou `id-MLDSA65-Ed25519`
  \item \textbf{Validation stricte} : Rejet de tout autre algorithme classique
  \item \textbf{Taillage} : Public key 1952 bytes, Secret key 2560 bytes
\end{itemize}

\section{CompositeJwk (Clés Hybrides)}

\subsection{Architecture Composite}

\begin{lstlisting}[language=Rust]
pub struct CompositeJwk {
    pub alg_id: CompositeAlgId,
    pub traditional_public_key: TraditionalJwk,
    pub pq_public_key: PostQuantumJwk,
}

pub enum CompositeAlgId {
    IdMldsa44Ed25519,  // DER: 060B6086480186FA6B5008013E
    IdMldsa65Ed25519,  // DER: 060B6086480186FA6B50080147
}
\end{lstlisting}

\subsection{Sérialisation JWK Composite}

\begin{lstlisting}[language=json]
{
  "algId": "id-MLDSA65-Ed25519",
  "traditionalPublicKey": {
    "kty": "OKP",
    "crv": "Ed25519",
    "x": "..."
  },
  "pqPublicKey": {
    "kty": "Akp",
    "crv": "ML-DSA-65",
    "x": "..."
  }
}
\end{lstlisting}

\subsection{Avantages de l'Architecture Composite}

\begin{enumerate}
  \item \textbf{Résistance Quantique} : ML-DSA-65 garantit sécurité post-quantique
  \item \textbf{Compatibilité} : Ed25519 assure compatibilité avec systèmes existants
  \item \textbf{Flexibilité} : Permet migration progressive vers PQC
  \item \textbf{Sécurité Améliorée} : Double validation augmente confiance
\end{enumerate}

% ============================================================================
\chapter{Cryptographie Post-Quantique}
% ============================================================================

\section{ML-DSA-65 (Module-Lattice Digital Signature)}

\subsection{Conformité et Standardisation}

\begin{itemize}
  \item \textbf{Norme} : NIST FIPS 204 (Approuvé pour usage gouvernemental)
  \item \textbf{Base} : Dilithium standardisé par NIST
  \item \textbf{Implémentation} : Open Quantum Safe (liboqs-rust)
  \item \textbf{Date d'approbation} : August 2024
\end{itemize}

\subsection{Spécifications Techniques}

\begin{table}[H]
\centering
\begin{tabular}{|l|r|r|}
\hline
\textbf{Paramètre} & \textbf{ML-DSA-44} & \textbf{ML-DSA-65} \\
\hline
Public Key Size & 1232 bytes & 1952 bytes \\
\hline
Secret Key Size & 2400 bytes & 2560 bytes \\
\hline
Signature Size & 2420 bytes & 3309 bytes \\
\hline
Security Level & Level 2 (128-bit) & Level 3 (192-bit) \\
\hline
Performance & $\sim$2-3 ms & $\sim$3-5 ms \\
\hline
\end{tabular}
\caption{Spécifications ML-DSA}
\end{table}

\subsection{Génération et Utilisation}

\begin{lstlisting}[language=Rust]
use pqc_ml_dsa::PqcSigner;

// Générer nouvelle paire de clés
let signer = PqcSigner::generate()?;

// Signer un message
let message = b"Document to sign";
let signature = signer.sign(message)?;

// Vérifier une signature
let is_valid = signer.verify(message, &signature)?;

// Récupérer depuis clés stockées
let signer = PqcSigner::from_secret_key(
    secret_key_bytes,
    public_key_bytes
)?;
\end{lstlisting}

\subsection{Résistance Quantique}

ML-DSA-65 résiste aux attaques quantiques via :
\begin{itemize}
  \item \textbf{Problème Lattice} : Très difficile même pour ordinateurs quantiques
  \item \textbf{Espace de clé énorme} : $2^{256}$ combinaisons possibles
  \item \textbf{Pas d'algorithme connu} : Pas d'équivalent quantique efficace du GNFS
\end{itemize}

\section{ML-KEM-768 (Key Encapsulation Mechanism)}

\subsection{Conformité et Standardisation}

\begin{itemize}
  \item \textbf{Norme} : NIST FIPS 203 (Kyber standardisé)
  \item \textbf{Base} : Kyber post-quantum KEM
  \item \textbf{Implémentation} : liboqs-rust
  \item \textbf{Type} : Key Encapsulation Mechanism (pas signature)
\end{itemize}

\subsection{Spécifications Techniques}

\begin{table}[H]
\centering
\begin{tabular}{|l|r|}
\hline
\textbf{Paramètre} & \textbf{Valeur} \\
\hline
Public Key Size & 1184 bytes \\
\hline
Secret Key Size & 2400 bytes \\
\hline
Ciphertext Size & 1088 bytes \\
\hline
Shared Secret Size & 32 bytes \\
\hline
Security Level & 192-bit (Level 3) \\
\hline
\end{tabular}
\caption{Spécifications ML-KEM-768}
\end{table}

\subsection{ML-KEM + AES-256-GCM (AEAD)}

Le mécanisme d'encapsulation ML-KEM-768 est combiné avec le chiffrement AES-256-GCM :

\begin{table}[H]
\centering
\begin{tabular}{|l|l|}
\hline
\textbf{ML-KEM-768} & \textbf{Génère shared secret 32 bytes} \\
\hline
\textbf{AES-256-GCM} & \textbf{Chiffre message avec authentification} \\
\hline
\textbf{Nonce/IV} & \textbf{12 bytes aléatoires via OsRng} \\
\hline
\end{tabular}
\end{table}

\subsection{Encapsulation dans les Credentials}

\begin{lstlisting}[language=Rust]
pub struct EncryptedMessage {
    pub kem_ciphertext: Vec<u8>,     // 1088 bytes (ML-KEM)
    pub encrypted_data: Vec<u8>,     // AES-GCM ciphertext
    pub nonce: Vec<u8>,              // 12 bytes
}

// Encapsulation
let km = MlKem768::encapsulate(issuer_public_key)?;
let encrypted = MlKem768::encrypt_message(
    &km.shared_secret,
    credential_bytes
)?;

// Décapsulation
let shared_secret = MlKem768::decapsulate(
    holder_secret_key,
    &encrypted.kem_ciphertext
)?;
let plaintext = MlKem768::decrypt_message(
    &shared_secret,
    &encrypted
)?;
\end{lstlisting}

% ============================================================================
\chapter{Signatures Hybrides (Weak Non-Separability)}
% ============================================================================

\section{Concept et Motivation}

\subsection{Le Problème}

L'adoption de la cryptographie post-quantique pose une question :

\begin{quote}
\textit{Comment transitionner de manière sécurisée des systèmes classiques vers post-quantiques ?}
\end{quote}

Deux approches :

\begin{enumerate}
  \item \textbf{Migration pure} : Abandonner Ed25519 pour ML-DSA-65 (risque de vulnérabilités non découvertes)
  \item \textbf{Approche hybride} : Combiner les deux (sécurité renforcée)
\end{enumerate}

\subsection{Signatures Hybrides}

Combinent ML-DSA-65 (post-quantum) et Ed25519 (classique) en une \textbf{signature composite} où chaque signature valide l'autre.

\section{Architecture HybridSigner}

\subsection{HybridKeyPair}

\begin{lstlisting}[language=Rust]
pub struct HybridKeyPair {
    pub pq_public_key: Vec<u8>,        // ML-DSA-65 public (1952 bytes)
    pub pq_secret_key: Vec<u8>,        // ML-DSA-65 secret (2560 bytes)
    pub classical_public_key: Vec<u8>, // Ed25519 public (32 bytes)
    pub classical_secret_key: Vec<u8>, // Ed25519 secret (32 bytes)
    pub key_id: String,                // blake3(pq_pub || ed_pub)
    pub created_at: u64,               // Timestamp création
}
\end{lstlisting}

\subsection{Génération de Clés Hybrides}

\begin{lstlisting}[language=Rust]
let signer = HybridSigner::generate()?;
// Generates:
// - ML-DSA-65 key pair (3.5 KB)
// - Ed25519 key pair (64 bytes)
// - Derived key_id from both keys

let composite_jwk = signer.export_composite_jwk();
// Exporte clés publiques compatibles IOTA
\end{lstlisting}

\section{Processus de Signature Composite}

\subsection{Weak Non-Separability (WNS)}

Le concept clé : \textbf{Empêcher l'attaque de stripping} où un adversaire supprime la signature PQC faible pour garder seulement la classique.

Mécanisme :
\begin{itemize}
  \item Lier PQ signature et classical signature cryptographiquement
  \item Inclure entropy aléatoire pour éviter déterminisme
  \item Valider binding hash pour garantir intégrité du lien
\end{itemize}

\subsection{Étapes du Processus}

\subsubsection{1. Calcul du Hash du Document}

\begin{lstlisting}[language=Rust]
let doc_hash = blake3::hash(message);
let doc_hash_bytes = doc_hash.as_bytes();  // 32 bytes
\end{lstlisting}

\subsubsection{2. Génération d'Entropie}

\begin{lstlisting}[language=Rust]
let entropy = self.generate_entropy();  // 16 bytes random
\end{lstlisting}

L'entropie est : 
\begin{itemize}
  \item Générée de manière sécurisée via \texttt{OsRng}
  \item Incluse dans les deux signatures (binding)
  \item Vérifiée lors de la validation
\end{itemize}

\subsubsection{3. Signature Post-Quantique}

\begin{lstlisting}[language=Rust]
let context = "QSDID-HYBRID-v1" 
              || doc_hash_bytes 
              || entropy 
              || classical_verifying_key.as_bytes()
              || pq_signer.public_key_bytes();

let pq_signature = pq_signer.sign(&context)?;
// Result: 3309 bytes
\end{lstlisting}

La signature PQC signe le \textbf{contexte lié} qui inclut :
\begin{itemize}
  \item Version du protocole
  \item Hash du document
  \item Entropy aléatoire
  \item Clés publiques (binding)
\end{itemize}

\subsubsection{4. Signature Classique}

\begin{lstlisting}[language=Rust]
let classical_msg = "CLASSIC-WRAPPER-v1"
                    || doc_hash_bytes
                    || pq_signature        // BINDING!
                    || entropy
                    || classical_verifying_key.as_bytes();

let classical_sig = classical_signer.sign(&classical_msg);
// Result: 64 bytes
\end{lstlisting}

La signature classique signe le message qui \textbf{inclut la signature PQC} : création du lien binding.

\subsubsection{5. Binding Hash}

\begin{lstlisting}[language=Rust]
let mut hasher = blake3::Hasher::new();
hasher.update(&pq_signature);
hasher.update(classical_signature.to_bytes().as_ref());
let binding_hash = hasher.finalize().into();  // 32 bytes
\end{lstlisting}

\subsection{CompositeSignature Résultante}

\begin{lstlisting}[language=Rust]
pub struct CompositeSignature {
    pub classical_signature: Vec<u8>,    // 64 bytes (Ed25519)
    pub pq_signature: Vec<u8>,           // 3309 bytes (ML-DSA-65)
    pub algorithm: CompositeAlgId::IdMldsa65Ed25519,
    pub timestamp: u64,                  // Unix epoch
    pub binding_hash: [u8; 32],          // Blake3(PQ || Classic)
    pub document_hash: [u8; 32],         // Blake3(message)
    pub entropy: [u8; 16],               // Random context binding
}

// Total size: ~3.5 KB
\end{lstlisting}

\section{Vérification de Signature Composite}

\subsection{Processus de Vérification}

\begin{enumerate}
  \item \textbf{Extraire contexte} : Reconstruire le contexte PQ original
  \item \textbf{Vérifier Ed25519} : Valider signature classique sur classical\_msg
  \item \textbf{Vérifier ML-DSA-65} : Valider signature PQC sur contexte reconstructed
  \item \textbf{Vérifier binding hash} : \texttt{Blake3(pq\_sig || classic\_sig) == binding\_hash}
  \item \textbf{Vérifier document hash} : \texttt{Blake3(message) == stored\_doc\_hash}
  \item \textbf{Vérifier timestamp} : Intervalle valide (anti-replay)
\end{enumerate}

\subsection{Conditions de Rejet}

\begin{table}[H]
\centering
\begin{tabular}{|l|l|}
\hline
\textbf{Erreur} & \textbf{Raison} \\
\hline
\texttt{WnsViolation} & Signatures non cryptographiquement liées \\
\hline
\texttt{HashMismatch} & Document altéré detecté \\
\hline
\texttt{TimestampValidationFailed} & Replay attack détecté \\
\hline
\texttt{AlgorithmMismatch} & Algo attendu $\neq$ algo reçu \\
\hline
\texttt{EntropyMismatch} & Context modifié \\
\hline
\texttt{ClassicalVerificationFailed} & Signature Ed25519 invalide \\
\hline
\texttt{PqVerificationFailed} & Signature ML-DSA-65 invalide \\
\hline
\end{tabular}
\caption{Conditions de rejet de signature composite}
\end{table}

% ============================================================================
\chapter{Stockage et Vérification}
% ============================================================================

\section{StorageClient (IPFS/Blockchain Integration)}

\subsection{Architecture}

Le \texttt{StorageClient} gère la persistance décentralisée des credentials :
\begin{itemize}
  \item Stockage sur IPFS (contenu adressé)
  \item Enregistrement sur blockchain (anchoring)
  \item Vérification de signatures distribuée
\end{itemize}

\subsection{Endpoints REST}

\subsubsection{POST /api/v1/store}

Enregistre un credential signé :

\begin{lstlisting}[language=json]
Request Body:
{
  "claims": {...},
  "metadata": {
    "type": "PersonalIdentity",
    "issuer": "did:iota:main:xyz...",
    "holder": "did:iota:main:abc...",
    "credentialId": "uuid",
    "issued_at": "2026-04-16T10:00:00Z",
    "ai_result": {
      "fraud_score": 0.05,
      "heatmap_hash": "blake3hash..."
    }
  },
  "signature": "base64_encoded_composite_signature"
}

Response:
{
  "success": true,
  "credentialId": "uuid",
  "rootCid": "QmXxxx..."  // IPFS CID
}
\end{lstlisting}

\subsubsection{GET /api/v1/retrieve/{rootCid}}

Récupère et valide un credential :

\begin{lstlisting}[language=json]
Query Parameters:
?issuerPubKey=hex_encoded_public_key

Response:
{
  "success": true,
  "claims": {...},
  "metadata": {...},
  "image": "base64_image",
  "signatureValid": true
}
\end{lstlisting}

\subsubsection{GET /health}

Heartbeat du service :

\begin{lstlisting}
Response: HTTP 200 OK
\end{lstlisting}

\section{PQC Verifier}

\subsection{Trait JwsVerifier}

L'interface \texttt{JwsVerifier} implémente la vérification généralisée :

\begin{lstlisting}[language=Rust]
pub trait JwsVerifier {
    fn verify_signature(
        &self,
        message: &[u8],
        signature: &[u8],
        public_key: &Jwk
    ) -> Result<bool>;
}
\end{lstlisting}

\subsection{Intégration Pipeline JOSE}

\begin{enumerate}
  \item Recevoir JWS avec algorithm header
  \item Extraire public key (Jwk)
  \item Sélectionner verifier approprié (classique ou PQC)
  \item Valider signature
  \item Retourner résultat
\end{enumerate}

\subsection{Support d'Algorithmes}

\begin{itemize}
  \item \textbf{Classiques} : ES256, EdDSA, RS256, etc.
  \item \textbf{Post-Quantiques} : id-MLDSA44-Ed25519, id-MLDSA65-Ed25519
  \item \textbf{Détection automatique} : Basée sur algorithm header
\end{itemize}

% ============================================================================
\chapter{Flux Complet d'un Credential}
% ============================================================================

\section{Flux End-to-End}

\subsection{Phase 1: Création d'Identité}

\begin{enumerate}
  \item Émetteur génère DID : \texttt{did:iota:main:hash(public\_keys)}
  \item Génère HybridKeyPair (ML-DSA-65 + Ed25519)
  \item Crée CustomCompositeJwk pour exportation
  \item Stocke clés privées dans HSM/Vault sécurisé
  \item Publie clés publiques sur blockchain
\end{enumerate}

\subsection{Phase 2: Création Credential}

\begin{enumerate}
  \item Émetteur collecte données holder
  \item Crée CredentialData :
  \begin{itemize}
    \item claims : données personnelles
    \item metadata : émetteur/holder DID, uuid
    \item ai\_result : fraude score
    \item image : photo document
  \end{itemize}
  \item Valide intégrité des données
\end{enumerate}

\subsection{Phase 3: Signature Composite}

\begin{enumerate}
  \item Blake3(credential) → doc\_hash
  \item Random(16) → entropy
  \item ML-DSA-65.sign(context) → pq\_signature (3309 bytes)
  \item Ed25519.sign(classical\_msg avec pq\_sig) → 64 bytes
  \item Blake3(pq\_sig || classic\_sig) → binding\_hash
  \item Composé : CompositeSignature (~3.5 KB)
\end{enumerate}

\subsection{Phase 4: Stockage Distribué}

\begin{enumerate}
  \item Encodage JSON + base64
  \item POST → StorageClient.store\_credential()
  \item Blockchain anchoring (transaction)
  \item IPFS storage (rootCid = Qm...)
  \item Retour credentialId et rootCid au holder
\end{enumerate}

\subsection{Phase 5: Vérification}

Holder/Verifier récupère credential :

\begin{enumerate}
  \item GET /api/v1/retrieve/{rootCid}?issuerPubKey={hex}
  \item Extraction issuer\_did
  \item Récupération public key depuis ledger
  \item Validation 6 conditions :
  \begin{enumerate}
    \item Ed25519 signature valide ?
    \item ML-DSA-65 signature valide ?
    \item Binding hash correspondant ?
    \item Document hash correspond ?
    \item Timestamp dans fenêtre valide ?
    \item Entropy consistent ?
  \end{enumerate}
  \item Résultat : credential accepté ✓ ou rejeté ✗
\end{enumerate}

\section{Diagramme de Flux}

\begin{lstlisting}
┌──────────────────────────────────────────────────────┐
│ 1. CREATION D'IDENTITÉ                               │
├──────────────────────────────────────────────────────┤
│  DID ← hash(public_keys)                             │
│  HybridKeyPair: ML-DSA-65 + Ed25519                 │
│  CustomCompositeJwk pour exportation                 │
└────────────────────────┬─────────────────────────────┘
                         │
                         ▼
┌──────────────────────────────────────────────────────┐
│ 2. CREATION CREDENTIAL                               │
├──────────────────────────────────────────────────────┤
│  CredentialData: claims, metadata, ai_result, image │
└────────────────────────┬─────────────────────────────┘
                         │
                         ▼
┌──────────────────────────────────────────────────────┐
│ 3. SIGNATURE COMPOSITE                               │
├──────────────────────────────────────────────────────┤
│  ML-DSA-65.sign(context) → 3309 bytes                │
│  Ed25519.sign(classical_msg) → 64 bytes              │
│  Blake3(pq_sig || classic_sig) → binding_hash        │
│  Total: ~3.5 KB                                      │
└────────────────────────┬─────────────────────────────┘
                         │
                         ▼
┌──────────────────────────────────────────────────────┐
│ 4. STOCKAGE DISTRIBUÉ                                │
├──────────────────────────────────────────────────────┤
│  StorageClient.store_credential()                    │
│  Blockchain anchoring + IPFS storage                 │
│  Response: rootCid = Qm...                           │
└────────────────────────┬─────────────────────────────┘
                         │
                         ▼
┌──────────────────────────────────────────────────────┐
│ 5. VERIFICATION (Holder/Verifier)                    │
├──────────────────────────────────────────────────────┤
│  GET /api/v1/retrieve/{rootCid}                      │
│  Validate 6 conditions:                              │
│    ✓ Ed25519 valid?                                  │
│    ✓ ML-DSA-65 valid?                                │
│    ✓ Binding hash match?                             │
│    ✓ Document hash match?                            │
│    ✓ Timestamp valid?                                │
│    ✓ Entropy consistent?                             │
│  Result: ✓ Accepted or ✗ Rejected                    │
└──────────────────────────────────────────────────────┘
\end{lstlisting}

% ============================================================================
\chapter{Sécurité et Protections}
% ============================================================================

\section{Matrice de Sécurité}

\begin{table}[H]
\centering
\begin{tabular}{|l|l|l|}
\hline
\textbf{Aspect} & \textbf{Protection} & \textbf{Mécanisme} \\
\hline
\textbf{Confidentialité} & ML-KEM-768 + AES-256-GCM & Chiffrement hybride PQC \\
\hline
\textbf{Intégrité} & Blake3 (double-hash) & Doc hash + signature hash \\
\hline
\textbf{Authenticité} & CompositeSignature & PQC + Classic binding \\
\hline
\textbf{Non-Répudiation} & Signatures digitales & Clés publiques vérifiables \\
\hline
\textbf{Résistance Quantique} & ML-DSA-65 (FIPS 204) & NIST standardisé \\
\hline
\textbf{Attaque Stripping} & WNS (Weak Non-Separability) & Contexte + entropie lié \\
\hline
\textbf{Replay Attack} & Timestamp validation & Fenêtre temporelle stricte \\
\hline
\textbf{Tampering} & Document hash & Detection modification \\
\hline
\textbf{Clés Compromises} & Revocation registry & Credential registry \\
\hline
\end{tabular}
\caption{Matrice de sécurité QSDID}
\end{table}

\section{Attaques Mitigées}

\subsection{Attaque Quantique}

\begin{itemize}
  \item \textbf{ML-DSA-65} : Fondé sur problème lattice difficile même quantiquement
  \item \textbf{Hardness Assumption} : LWE (Learning With Errors) pas d'algorithme poly-temps
  \item \textbf{Sécurité} : 128-bit classique = 192-bit lattice
\end{itemize}

\subsection{Attaque de Stripping}

\begin{itemize}
  \item \textbf{Problème} : Adversaire supprime PQC, garde seulement classique
  \item \textbf{Solution} : WNS - inclure PQ\_sig dans classical\_msg
  \item \textbf{Validation} : Vérifier binding hash = Blake3(PQ || Classic)
\end{itemize}

\subsection{Attaque Replay}

\begin{itemize}
  \item \textbf{Problème} : Replay old credentials
  \item \textbf{Solution} : Timestamp + nonce + entropy
  \item \textbf{Validation} : Temps dans $\pm 5$ minutes (configurable)
\end{itemize}

\subsection{Attaque Man-in-the-Middle}

\begin{itemize}
  \item \textbf{Problème} : Interception et modification
  \item \textbf{Solution} : Blake3(document) + dual signatures
  \item \textbf{Détection} : HashMismatch error
\end{itemize}

\section{Gestion des Clés}

\subsection{Cycles de Vie}

\begin{enumerate}
  \item \textbf{Génération} : OsRng cryptographiquement sécurisé
  \item \textbf{Stockage} : HSM/Vault (clés privées jamais en RAM)
  \item \textbf{Backup} : Chiffrement + distribution géographique
  \item \textbf{Rotation} : Nouvelle paire, revocation ancienne
  \item \textbf{Révocation} : Enregistrement dans registry
\end{enumerate}

\subsection{Memory Safety}

\begin{itemize}
  \item \textbf{Rust} : Pas d'accès mémoire après-libération
  \item \textbf{zeroize} : Effacement sécurisé de clés en RAM
  \item \textbf{OsRng} : RNG système hautement entropique
\end{itemize}

% ============================================================================
\chapter{Dépendances Cryptographiques}
% ============================================================================

\section{Stack Cryptographique}

\begin{lstlisting}
Workspace dependencies:
├── oqs (0.11)              # liboqs-rust (OpenQuantumSafe)
│   └─ ML-DSA-65 (FIPS 204)
│   └─ ML-KEM-768 (FIPS 203)
├── iota-crypto (0.23.2)    # Ed25519, SHA
├── ed25519-dalek (2.1)     # Signatures Ed25519
├── aes-gcm (0.10)          # Chiffrement authentifié
├── blake3 (1.5)            # Hachage cryptographique
├── serde_json              # Sérialisation JSON
├── bls12_381_plus          # Courbes elliptiques
├── zeroize                 # Effacement mémoire
├── anyhow, thiserror       # Gestion erreurs
├── hex, base64             # Encodages
└── rand (0.8.5)            # Entropie CSPRNG
\end{lstlisting}

\subsection{Chaîne de Confiance}

\begin{enumerate}
  \item \textbf{NIST} approuve ML-DSA-65 et ML-KEM-768
  \item \textbf{OpenQuantumSafe} implémente avec rigor
  \item \textbf{liboqs-rust} binding sécurisé
  \item \textbf{iota-crypto} utilitaires validés
  \item \textbf{Projet QSDID} intégration et validation
\end{enumerate}

% ============================================================================
\chapter{Architecture Globale}
% ============================================================================

\section{Diagramme d'Architecture}

\begin{lstlisting}
┌────────────────────── QSDID Platform ──────────────┐
│                                                     │
│  ┌─── Identity Layer ───────────────────────────┐ │
│  │ DID + CredentialData + Metadata              │ │
│  │ (Decentralized + Verifiable)                 │ │
│  └────────────────────────────────────────────┘ │
│                     │                             │
│  ┌─── Crypto Layer ──────────────────────────┐ │ │
│  │                                            │ │ │
│  │  ┌─ Classical ────┐  ┌─ Post-Quantum ──┐ │ │ │
│  │  │ Ed25519 (64B)  │  │ ML-DSA-65      │ │ │ │
│  │  │ K256/P256      │  │ (3309B)        │ │ │ │
│  │  │                │  │ ML-KEM-768     │ │ │ │
│  │  └────────────────┘  └────────────────┘ │ │ │
│  │        ▲                      ▲          │ │ │
│  │        └──────────┬───────────┘          │ │ │
│  │                ▼                         │ │ │
│  │          Hybrid Signatures               │ │ │
│  │      (WNS Binding + Entropy)            │ │ │
│  │                                          │ │ │
│  └──────────────────────────────────────────┘ │ │
│                     │                          │
│  ┌─── JWK Encapsulation ───────────────────┐ │ │
│  │                                          │ │ │
│  │  TraditionalJwk  CompositeJwk  PQCJwk   │ │ │
│  │  (RSA,EC,OKP)    (Hybrid)      (AKP)    │ │ │
│  │                                          │ │ │
│  └──────────────────────────────────────────┘ │ │
│                     │                          │
│  ┌─── Storage Layer ────────────────────────┐ │ │
│  │ IPFS/Blockchain (rootCid) + Registry     │ │ │
│  │ StorageClient REST API                   │ │ │
│  └──────────────────────────────────────────┘ │ │
│                                                │
└────────────────────────────────────────────────┘
\end{lstlisting}

\section{Résumé Exécutif}

La plateforme QSDID implémente une architecture d'identité décentralisée hautement sécurisée :

\begin{enumerate}
  \item \textbf{Identités Décentralisées} : DIDs vérifiables, sans autorité centrale
  \item \textbf{Credentials Signés} : Signatures hybrides avec double validation
  \item \textbf{Résistance Quantique} : ML-DSA-65 et ML-KEM-768 (NIST FIPS 204/203)
  \item \textbf{Encapsulation Clés} : JWK, CompositeJwk, PostQuantumJwk
  \item \textbf{Stockage Distribué} : IPFS + Blockchain anchoring
  \item \textbf{Vérification Robuste} : Six étapes de validation de signature
  \item \textbf{Protection Multi-niveaux} : Confidentialité, intégrité, authenticité, non-répudiation
\end{enumerate}

% ============================================================================
\appendix

\chapter{Glossaire}

\begin{description}
  \item[DID] Decentralized Identifier (Identifiant Décentralisé)
  \item[JWK] JSON Web Key (Clé Web JSON)
  \item[JWS] JSON Web Signature (Signature Web JSON)
  \item[JWT] JSON Web Token (Token Web JSON)
  \item[JOSE] JSON Object Signing and Encryption
  \item[ML-DSA-65] Module-Lattice Digital Signature (NIST FIPS 204)
  \item[ML-KEM-768] Module-Lattice Key Encapsulation Mechanism (NIST FIPS 203)
  \item[PQC] Post-Quantum Cryptography
  \item[WNS] Weak Non-Separability
  \item[AEAD] Authenticated Encryption with Associated Data
  \item[OKP] Octet string key pairs
  \item[AKP] Asymmetric Key Pair (post-quantum)
  \item[HSM] Hardware Security Module
  \item[IPFS] InterPlanetary File System
  \item[CID] Content Identifier
  \item[CSPRNG] Cryptographically Secure Pseudo-Random Number Generator
\end{description}

\chapter{Ressources}

\section{Standards}

\begin{itemize}
  \item NIST FIPS 204 : Module-Lattice-Based Digital Signature Standard
  \item NIST FIPS 203 : Module-Lattice-Based Key-Encapsulation Mechanism
  \item W3C DID Core 1.0 : https://www.w3.org/TR/did-core/
  \item RFC 7517 : JSON Web Key (JWK)
  \item RFC 7515 : JSON Web Signature (JWS)
  \item RFC 7519 : JSON Web Token (JWT)
\end{itemize}

\section{Références Utiles}

\begin{itemize}
  \item Open Quantum Safe : https://openquantumsafe.org/
  \item IOTA Identity Framework : https://docs.iota.org/iota-identity
  \item Rust Cryptography : https://www.rust-lang.org/what/wg-crypto/
  \item Post-Quantum Cryptography : https://csrc.nist.gov/projects/post-quantum-cryptography/
\end{itemize}

\end{document}
